\documentclass{article}

\usepackage{amsmath,amssymb}
\usepackage{xurl}
\usepackage[hidelinks]{hyperref}
\setlength{\emergencystretch}{2em}

% Shared commands for notes in docs/paper-gaps/.

\DeclareUnicodeCharacter{2080}{\ifmmode{}_{0}\else\textsubscript{0}\fi}
\DeclareUnicodeCharacter{2081}{\ifmmode{}_{1}\else\textsubscript{1}\fi}
\DeclareUnicodeCharacter{2082}{\ifmmode{}_{2}\else\textsubscript{2}\fi}
\DeclareUnicodeCharacter{2099}{\ifmmode{}_{n}\else\textsubscript{n}\fi}

\newcommand{\Lean}{\textsc{Lean}}
\ExplSyntaxOn
\NewDocumentCommand{\leanid}{m}{
  \ifmmode
    \textsf{\detokenize{#1}}
  \else
    \group_begin:
    \urlstyle{sf}
    \tl_set:Nn \l_tmpa_tl {#1}
    \tl_replace_all:Nnn \l_tmpa_tl {\_} {_}
    \exp_args:NV \path \l_tmpa_tl
    \group_end:
  \fi
}
\ExplSyntaxOff
\providecommand{\tr}{\operatorname{tr}}
\newcommand{\tnleanIssueBase}{https://github.com/LionSR/TNLean/issues/}
\newcommand{\tnleanPrBase}{https://github.com/LionSR/TNLean/pull/}
\newcommand{\tnleanDocBase}{https://sirui-lu.com/TNLean/docs/}
\newcommand{\ghissue}[1]{\href{\tnleanIssueBase#1}{GitHub issue \##1}}
\newcommand{\ghpr}[1]{\href{\tnleanPrBase#1}{PR \##1}}
\newcommand{\doclink}[1]{\href{\tnleanDocBase#1.html}{\path{#1}}}

% Dirac notation and the identity operator, as in blueprint/src/macros/common.tex.
% \providecommand keeps note-local definitions authoritative.
\usepackage{dsfont}
\providecommand{\ket}[1]{|#1\rangle}
\providecommand{\bra}[1]{\langle #1|}
\providecommand{\braket}[2]{\langle #1 | #2 \rangle}
\providecommand{\ketbra}[2]{|#1\rangle\!\langle #2|}
\providecommand{\Id}{\mathds{1}}
\providecommand{\one}{\mathds{1}}
\providecommand{\C}{\mathbb{C}}
\providecommand{\MN}[1]{M_{#1}(\C)}
\providecommand{\MD}{M_D(\C)}

% External referencing for paper sources.  The build helper writes sanitized
% auxiliary files under build/paper-aux/ and excludes duplicated source labels.
\usepackage{xr}

% Import a generated paper auxiliary file with a caller-supplied label prefix.
% The candidate paths support builds from docs/paper-gaps/, the repository root,
% and build/paper-aux/ itself.
\newcommand{\importpaperaux}[2]{%
  \IfFileExists{../../build/paper-aux/#2.aux}{%
    \externaldocument[#1]{../../build/paper-aux/#2}%
  }{%
    \IfFileExists{build/paper-aux/#2.aux}{%
      \externaldocument[#1]{build/paper-aux/#2}%
    }{%
      \IfFileExists{#2.aux}{%
        \externaldocument[#1]{#2}%
      }{}%
    }%
  }%
}

% General external reference macros
\newcommand{\paperref}[2]{\ref{#1-#2}}
\newcommand{\papereqref}[2]{(\ref{#1-#2})}

% CPSV16 source (1606.00608 / MPDO-22-12-17-2.tex) external document and shortcuts
\importpaperaux{cpsv16-}{1606.00608/MPDO-22-12-17-2-tnlean}
\newcommand{\cpsvref}[1]{\paperref{cpsv16}{#1}}
\newcommand{\cpsveqref}[1]{\papereqref{cpsv16}{#1}}

% Verdict marker: \gapnote{<kind>}{<status>}. Kind is the policy
% classification (clarification, local-correction, scope-restriction,
% unfaithful, false-source, open-gap); status is open, wip (elimination
% underway), resolved, or historical. Machine-read by the site index;
% prints a verdict line.
\newcommand{\gapnote}[2]{\par\noindent\textbf{Verdict:} #1 (#2).\par}


\title{The Nonzero-Coefficient Convention for CPSV16 Canonical Forms}
\author{}
\date{2026-08-23}

\begin{document}
\maketitle
\gapnote{local-correction}{resolved}

\section{The convention}

Cirac--P\'erez-Garc\'ia--Schuch--Verstraete construct the canonical form by
listing the invariant blocks of the completely positive map and normalizing
each block separately \cite[lines 214--225]{Cirac2016MPDO_arXiv}. The blocks
produced by the construction are therefore precisely the blocks which occur in
the tensor. The displayed definition
\cite[lines 238--244]{Cirac2016MPDO_arXiv},
\begin{align}
    A^i=\bigoplus_{k=1}^r \mu_k A_k^i,
    \tag{\path{eq:II_CF1}}
\end{align}
does not repeat that every $\mu_k$ is nonzero, but the subsequent
normalization $|\mu_k|\le1$ with at least one coefficient equal to one
\cite[line 246]{Cirac2016MPDO_arXiv} presupposes it: a summand with
coefficient zero contributes nothing to any positive-length matrix product
vector and is not a block of the decomposition.

The formalization reads the canonical form with every coefficient $\mu_k$
nonzero. This is part of the definition of canonical-form data.\footnote{The
field is \leanid{MPSTensor.CPSVCanonicalFormData.weights_ne_zero} in
\doclink{TNLean/MPS/CanonicalForm/Definitions.lean}.} A basis of normal
tensors in the sense of Definition~2.4 is then a family of normal tensors
whose matrix product vectors span those of $A$ at every positive length and
are eventually linearly independent; the blocks of a canonical form form
such a basis after identifying gauge-phase-equivalent blocks, and every
basis element occurs with nonzero coefficient.

\section{Consequences}

Under this convention the following source statements are formalized as
stated, each against the blueprint node named.

\begin{itemize}
    \item Proposition~2.7 \cite[lines 278--280 and 1135--1146]
          {Cirac2016MPDO_arXiv}: a family is a basis of normal tensors for a
          canonical form if and only if every representative is normal, every
          block is gauge-phase equivalent to a representative, and distinct
          representatives are not gauge-phase equivalent. Blueprint
          \path{thm:cpsv_bnt_characterization}.\footnote{Formal declaration:
          \leanid{MPSTensor.isCPSVBasisOfNormalTensors_iff_canonicalForm_covered_and_minimal}.}
    \item The uniqueness sentence at \cite[line 1148]{Cirac2016MPDO_arXiv}:
          two bases of normal tensors of the same positive-length matrix
          product vector family have a bijective gauge-phase matching, proved
          by applying Proposition~2.7 in both directions. It is formalized
          for basis-of-normal-tensors sector presentations, in which every
          basis element carries a nonzero coefficient. Blueprint
          \path{thm:cpsv_bnt_presentation_uniqueness}.\footnote{Formal
          declaration:
          \leanid{MPSTensor.IsBNTSectorPresentation.equiv_of_sameMPV₂Pos}.}
          The literal Definition-2.4 family predicate\footnote{Formal
          declaration: \leanid{MPSTensor.IsCPSVBasisOfNormalTensors}.} is
          retained as an interface and carries no uniqueness claim: for a
          one-block tensor $A_0$ with coefficient $1$, both $(A_0)$ and
          $(A_0,A_1)$ satisfy the family predicate, the second by assigning
          $A_1$ the coefficient zero at every length. That second family is
          exactly a zero-coefficient member, which the convention excludes
          from every basis presentation, so it is not a counterexample to
          the formalized uniqueness statement.
    \item Theorem~2.10, Corollary~2.11, and Corollary~A.6
          \cite[lines 349--361 and 1167--1199]{Cirac2016MPDO_arXiv}: the
          proportional and equal fundamental theorems and their unitary
          refinements. Blueprint \path{thm:cpgsv_multiblock_ft_source},
          \path{thm:cpgsv_equal_case_source},
          \path{cor:sector_bnt_proportional_unitary_sector_match}, and
          \path{cor:sector_bnt_equal_unitary_global_gauge}.
    \item Lemma~A.5 \cite[lines 1155--1163]{Cirac2016MPDO_arXiv}: equal
          finite-range power sums of nonzero lists give equal lengths and
          equal multisets; it is applied to the nonzero coefficients of a
          canonical form. Blueprint \path{thm:bounded_power_sum_multiset}.
    \item Corollary~3.12 \cite[lines 583--590]{Cirac2016MPDO_arXiv}: every
          block of a canonical-form renormalization fixed point has
          unit-modulus coefficient and idempotent transfer map, and the
          phase-class representatives have a joint residual isometry.
          Blueprint \path{thm:cpsv_weight_norm_one_and_block_rfp} and
          \path{thm:cpsv_phase_class_residual_isometries}.
\end{itemize}

Without the convention, a normal tensor with coefficient identically zero
could be adjoined to any basis without changing a positive-length matrix
product vector, and the cardinality conclusions of the uniqueness sentence,
Theorem~2.10, and Lemma~A.5 would fail. The convention is the source's own
reading, and it is recorded here so that the formal definition can be checked
against the displayed one.

\bibliographystyle{alpha}
\bibliography{references}

\end{document}
